\chapter{Wprowadzenie}\label{ch:wprowadzenie}

\section{Kontekst technologiczny}\label{sec:kontekst_technologiczny}

Dynamiczny rozwój technologii informatycznych oraz powszechna cyfryzacja procesów biznesowych i instytucjonalnych sprawiły, że cyberbezpieczeństwo stało się jednym z kluczowych filarów ciągłości działania współczesnych organizacji. Równolegle do postępu w obszarze mechanizmów defensywnych następuje ciągła ewolucja metod i wektorów ataków stosowanych przez ugrupowania cyberprzestępcze oraz zaawansowane grupy sponsorowane przez państwa (ang. \textit{Advanced Persistent Threat}, APT). Tradycyjny paradygmat ochrony stacji roboczych i serwerów, opierający się na klasycznych programach antywirusowych (ang. \textit{Endpoint Protection Platforms}, EPP) weryfikujących sygnatury plików statycznie zapisanych na wolumenie dyskowym, wykazuje obecnie daleko idącą nieskuteczność.

Współczesne rodziny złośliwego oprogramowania, w szczególności konie trojańskie zdalnego dostępu (ang. \textit{Remote Access Trojan}, RAT, np. Agent Tesla, AsyncRAT, Remcos) oraz wyspecjalizowane oprogramowanie kradnące tożsamość cyfrową i poświadczenia (ang. \textit{Infostealer}, np. RedLine, Vidar, LummaC2), wykorzystują złożone mechanizmy zaciemniania kodu (ang. \textit{obfuscation}) oraz unikania mechanizmów detekcji (ang. \textit{defense evasion}). Do fundamentalnych technik utrudniających statyczną analizę należą:
\begin{itemize}
    \item \textbf{Techniki bezplikowe oraz Living-off-the-Land (LotL / LOLBins):} wykorzystywanie wbudowanych narzędzi administracyjnych systemu operacyjnego (takich jak \code{PowerShell}, \code{WMI}, \code{certutil}, \code{mshta}) do pobierania, dekodowania i uruchamiania ładunków w sposób niewymagający zapisu binariów na nośniku pamięci masowej.
    \item \textbf{Wstrzykiwanie kodu do przestrzeni procesów:} zacieranie śladów poprzez alokację i egzekucję złośliwego kodu maszynowego w kontekście bezpieczeństwa zaufanych procesów systemowych (np. \file{explorer.exe}, \file{svchost.exe}) przy użyciu funkcji API jądra Windows (\code{VirtualAllocEx}, \code{WriteProcessMemory}, \code{CreateRemoteThread}).
    \item \textbf{Drążenie procesów (ang. \textit{Process Hollowing}):} technika polegająca na powołaniu legalnego procesu w stanie zawieszenia (\textit{suspended state}), odmapowaniu oryginalnej sekcji kodu (\code{NtUnmapViewOfSection}) i wstrzyknięciu w jej miejsce złośliwego pliku wykonywalnego przed wznowieniem wątku głównego.
    \item \textbf{Refleksywne ładowanie bibliotek (ang. \textit{Reflective DLL Injection}):} załadowanie biblioteki dynamicznej bezpośrednio z bufora pamięci RAM bez angażowania standardowego mechanizmu ładowania systemu (\code{ntdll!LdrLoadDll}) i bez rejestracji w oficjalnych strukturach modułów procesu (\code{PEB.Ldr}).
\end{itemize}

Powyższy stan rzeczy determinuje przesunięcie punktu ciężkości w operacjach bezpieczeństwa (ang. \textit{Security Operations Center}, SOC) oraz w informatyce śledczej i reagowaniu na incydenty (ang. \textit{Digital Forensics and Incident Response}, DFIR). Niezbędna staje się analiza ulotnych śladów binarnych w pamięci operacyjnej (RAM) oraz ciągła telemetria behawioralna jądra systemu operacyjnego. Ulotna pamięć RAM stanowi najbardziej wiarygodne źródło wiedzy o incydencie, gdyż przechowuje wstrzyknięte ładunki w postaci odszyfrowanej, struktury procesów, poświadczenia uwierzytelniające oraz aktywne deskryptory połączeń sieciowych.

\section{Problem badawczy i motywacja}\label{sec:problem_badawczy}

Współczesna inżynieria detekcji (ang. \textit{Detection Engineering}) zmaga się z dwoma krytycznymi wyzwaniami:
\begin{enumerate}
    \item \textbf{Zjawiskiem przeciążenia alertami (ang. \textit{alert fatigue}):} wynikającym z generowania nadmiernej liczby fałszywych alarmów (ang. \textit{False Positives}, FP) przez niedostatecznie zoptymalizowane lub zbyt ogólne reguły detekcyjne.
    \item \textbf{Niską odpornością wskaźników kompromitacji na modyfikacje:} zjawiskiem opisanym w koncepcji Piramidy Bólu Davida Bianco (ang. \textit{Pyramid of Pain}). Tradycyjne wskaźniki kompromitacji (ang. \textit{Indicators of Compromise}, IoC), takie jak skróty kryptograficzne plików (MD5, SHA256) czy adresy IP serwerów dowodzenia (ang. \textit{Command and Control}, C2), są trywialne do zneutralizowania przez adwersarza poprzez zwykłą rekompilację lub rotację infrastruktury.
\end{enumerate}

Aby skutecznie przeciwdziałać zaawansowanym zagrożeniom, konieczne jest tworzenie reguł detekcyjnych operujących na wyższych piętrach piramidy - wskaźnikach aktywności (ang. \textit{Indicators of Activity}, IoA) oraz technikach, taktykach i procedurach (ang. \textit{Tactics, Techniques, and Procedures}, TTP) skatalogowanych w matrycy MITRE ATT\&CK.

Wytworzenie precyzyjnych reguł wymaga jednak przeprowadzenia kontrolowanej infekcji w środowisku laboratoryjnym, pozyskania zrzutu pamięci operacyjnej oraz rejestracji szczegółowych dzienników zdarzeń (np. za pośrednictwem narzędzia Microsoft Sysmon). Następnie konieczne jest przeprowadzenie korelacji artefaktów z pamięci RAM (ujawnionych np. przez Volatility 3 lub MemProcFS) z telemetrią logów systemowych (parsowanych przy użyciu nowoczesnych narzędzi, takich jak Hayabusa czy Chainsaw).

Istotną luką w wielu opracowaniach laboratoryjnych jest brak spójnego, powtarzalnego podejścia, które w sposób zintegrowany łączy fazę inżynierii laboratoryjnej, poinfekcyjną analizę śledczą oraz bezpośrednią syntezę pozyskanych danych w postaci reguł YARA (dla pamięci RAM) oraz Sigma (dla dzienników zdarzeń). Niniejszy projekt stanowi odpowiedź na tę potrzebę.

\section{Cel i teza projektu}\label{sec:cele_projektu}

Głównym celem projektu jest \textbf{zaprojektowanie, wdrożenie i walidacja bezpiecznego, odizolowanego środowiska laboratoryjnego oraz opracowanie metodyki ekstrakcji i korelacji poinfekcyjnych artefaktów pamięci RAM i dzienników zdarzeń Windows w celu wytworzenia i weryfikacji wysokowiernych reguł detekcyjnych YARA i Sigma, a także skompilowania kompleksowej matrycy IoC/IoA.}

W ramach celu głównego wyznaczono następujące cele szczegółowe:
\begin{enumerate}
    \item \textbf{Opracowanie architektury i wdrożenie laboratorium badawczego:}
    Utworzenie odizolowanego segmentu wirtualnego, obejmującego stację ofiary (system Windows 10/11) ze specjalnie ukształtowanym profilem telemetrii Sysmon i dezaktywowanymi mechanizmami ochronnymi (Windows Defender, AMSI) oraz węzeł analityczny bazujący na systemie Linux (np. REMnux), wyposażony w usługi emulacji sieci (FakeNet-NG / INetSim).
    
    \item \textbf{Akwizycja i kontrolowana egzekucja materiału badawczego:}
    Bezpieczne pobranie z repozytoriów złośliwego oprogramowania (np. MalwareBazaar) próbek reprezentujących rodziny Infostealer oraz RAT, ich statyczna wstępna weryfikacja, detonacja w monitorowanym węźle oraz pozyskanie pełnego zrzutu pamięci operacyjnej z poziomu hiperwizora wraz z eksportem dzienników \file{.evtx}.
    
    \item \textbf{Informatyka śledcza pamięci ulotnej (Memory Forensics):}
    Identyfikacja ukrytych procesów, anomalii w strukturze drzewa procesów, zmapowanie nieprzypisanych pamięci VAD ze statusem \code{PAGE_EXECUTE_READWRITE} (detekcja iniekcji kodu) oraz ekstrakcja zdekodowanych ładunków binarnych bezpośrednio z RAM za pomocą narzędzi Volatility 3 oraz MemProcFS.
    
    \item \textbf{Analiza telemetrii systemowej i korelacja zdarzeń:}
    Zautomatyzowane przetworzenie logów systemowych za pomocą narzędzi Chainsaw oraz Hayabusa, rekonstrukcja łańcucha wykonania (drzewo procesów, parametry CLI, manipulacje rejestrem w celach persystencji, aktywność sieciowa) oraz powiązanie zachowań telemetrycznych z artefaktami pamięciowymi.
    
    \item \textbf{Inżynieria detekcji i weryfikacja empiryczna:}
    Zaprojektowanie autorskich reguł YARA ukierunkowanych na wykrywanie artefaktów pamięciowych w procesach oraz generycznych reguł Sigma wykrywających anomalie behawioralne w dziennikach systemowych, weryfikacja ich skuteczności oraz współczynnika fałszywych alarmów (FP) w czystym systemie referencyjnym, a także stworzenie matrycy wskaźników IoC oraz IoA.
\end{enumerate}

\section{Zakres projektu}\label{sec:zakres_projektu}

Zakres pracy obejmuje pięć głównych faz realizacyjnych, przedstawionych w tabeli~\ref{tab:fazy_projektu}.

\begin{table}[htbp]
\centering
\caption{Zestawienie faz realizacyjnych i komponentów projektu}
\label{tab:fazy_projektu}
\tablefont
\begin{tabularx}{\linewidth}{l L{4.2cm} Y}
\toprule
\hcell{Faza} & \hcell{Obszar prac} & \hcell{Główne narzędzia / technologie} \\
\midrule
I & Architektura i wdrożenie laboratorium & Hiperwizor (Type-2), izolowana sieć wewnętrzna (\textit{host-only}), FakeNet-NG, INetSim, PowerShell hardening/unhardening, Sysmon. \\
\addlinespace
II & Pozyskanie próbek i bezpieczna detonacja & Repozytorium MalwareBazaar, narzędzia statyczne (PEstudio, Floss), mechanizm snapshotów hiperwizora, akwizycja RAM i logów \file{.evtx}. \\
\addlinespace
III & Śledcza analiza pamięci RAM & Volatility 3, MemProcFS, wtyczki \code{malfind}, \code{pslist}, \code{netscan}, \code{vadinfo}, disasemblery (Ghidra, IDA). \\
\addlinespace
IV & Analiza i korelacja logów systemowych & Sysmon Modular, Chainsaw, Hayabusa, Windows Event Logs, parsery JSON/CSV, timeline analysis. \\
\addlinespace
V & Inżynieria detekcji i walidacja & Standard YARA (analiza pamięci procesu), standard Sigma (analityka logów), walidacja FP/TP, matryca IoC/IoA, MITRE ATT\&CK. \\
\bottomrule
\end{tabularx}
\end{table}

\subsection*{Ograniczenia i założenia projektowe}
W celu zachowania przejrzystości metodologicznej i powtarzalności eksperymentu przyjęto następujące założenia:
\begin{itemize}
    \item \textbf{Docelowy system operacyjny:} Eksperymenty przeprowadzane są w 64-bitowym środowisku Microsoft Windows (Windows 10/11), jako że stanowi ono dominujący cel w sektorze ataków na końcówki robocze.
    \item \textbf{Izolacja sieciowa:} Całkowity zakaz bezpośredniej komunikacji węzła infekcyjnego z publicznym Internetem; cała komunikacja C2 jest zamykana i emulowana lokalnie przez stację analityczną (FakeNet-NG / INetSim).
    \item \textbf{Faza ataku:} Badania koncentrują się na analizie poinfekcyjnej (\textit{post-execution / post-compromise}), kładąc główny nacisk na fazy wykonania (\textit{Execution}), utrzymania dostępu (\textit{Persistence}), unikania obrony (\textit{Defense Evasion}) oraz komunikacji (\textit{Command and Control}).
\end{itemize}

\section{Metodologia badawcza i warsztat technologiczny}\label{sec:metodologia}

Metodologia zastosowana w niniejszym projekcie opiera się na iteracyjnym cyklu inżynierii detekcji (ang. \textit{Detection Engineering Life Cycle}), który integruje analizę behawioralną, analizę statyczną artefaktów oraz formalne testowanie reguł detekcyjnych.

W toku realizacji projektu wykorzystywane są następujące grupy narzędzi:
\begin{itemize}
    \item \textbf{Środowisko wirtualizacji i symulacji sieci:} Wirtualne maszyny operujące w odciętej podsieci; węzeł analityczny przechwytuje zapytania DNS, protokoły HTTP/HTTPS, SMTP oraz surowe połączenia TCP, generując wiarygodne syntetyczne odpowiedzi serwera C2.
    \item \textbf{Telemetria systemowa (Sysmon):} Wykorzystanie rozbudowanego pliku konfiguracyjnego (m.in. na bazie standardu \textit{Sysmon-Modular}), pozwalającego na rejestrację m.in. powoływania wątków zdalnych (Event ID 8), uzyskiwania dostępu do procesów LSASS/innych (Event ID 10) oraz ładowania sterowników i bibliotek (Event ID 6 i 7).
    \item \textbf{Frameworki analizy pamięci (Volatility 3 / MemProcFS):} Zastosowanie technik eksploracji pamięci fizycznej i wirtualnej, omijających manipulacje adwersarza w tablicach jądra (\textit{DKOM - Direct Kernel Object Manipulation}) i pozwalających na ekstrakcję czystego kodu wykonywalnego z deskryptorów VAD.
    \item \textbf{Narzędzia szybkiego triage logów (Hayabusa / Chainsaw):} Przetwarzanie wielomegabajtowych plików \file{.evtx} z wykorzystaniem wbudowanych reguł Sigma w celu natychmiastowej rekonstrukcji osi czasu incydentu.
    \item \textbf{Formaty reguł detekcyjnych (YARA i Sigma):} Opracowanie reguł YARA w standardzie wersji 4.x ze szczególnym uwzględnieniem modułów \code{pe} i selekcji heksadecymalnej oraz reguł Sigma w formacie YAML zgodnych ze specyfikacją \textit{SigmaHQ}.
\end{itemize}
